\newglossaryentry{timple}
{
    name=timple,
    description={Pequeño instrumento de cinco cuerdas utilizado principalmente para el acompañamiento de canciones folkloricas en las Islas Canarias}
}

\newglossaryentry{feedback}
{
    name=feedback,
    description={Comentario o sugerencia de un usuario}
}

\newglossaryentry{notación musical}
{
    name=notación musical,
    description={Sistema de escritura utilizado para representar gráficamente una pieza musical}
}

\newglossaryentry{notación latina}
{
    name=notación latina,
    description={Notación musical para nombrar las notas musicales mediante las silabas do, re, mi, fa, sol, la y si}
}

\newglossaryentry{notación anglosajona}
{
    name=notación anglosajona,
    description={Notación musical para nombrar las notas musicales mediante las vocales C, D, E, F, G, A y B}
}

\newglossaryentry{escala musical}
{
    name=escala musical,
    description={Conjunto de sonidos ordenados para formar una secuencia de notas musicales}
}

\newglossaryentry{tónica}
{
    name=tónica,
    description={Nota cabecera de una escala}
}

\newglossaryentry{semitono}
{
    name=semitono,
    description={Distancia mínima entre dos notas}
}

\newglossaryentry{tono}
{
    name=tono,
    description={Equivalente a dos semitonos}
}


\newglossaryentry{octava}
{
    name=octava,
    description={Secuencia que forman las notas do, re, mi, fa, sol, la, si, do}
}

\newglossaryentry{acorde}
{
    name=acorde,
    description={Conjunto de tres o mas notas diferentes que en su reproducción constituyen una unidad armónica}
}

\newglossaryentry{cualidad}
{
    name=cualidad,
    description={Intervalo de semitonos que definen la construcción de un acorde}
}

\newglossaryentry{progresión}
{
    name=progresión,
    description={Sucesión de acordes}
}

\newglossaryentry{transposición}
{
    name=transposición,
    description={Incremento o decremento de los índices de las notas que conforman un acorde}
}

\newglossaryentry{afinación}
{
    name=afinación,
    description={En instrumentos de cuerda pulsada, acción de tensar una cuerda para obtener una nota}
}

\newglossaryentry{punteo}
{
    name=punteo,
    description={En instrumentos de cuerda pulsada, técnica que consiste en tocar piezas musicales tocando notas de forma individual o combinada}
}

\newglossaryentry{rasgueo}
{
    name=rasgueo,
    description={En instrumentos de cuerda pulsada, técnica que consiste en tocar distintas cuerdas utilizando uno o varios dedos, en forma de barrido}
}

\newglossaryentry{cuerda al aire}
{
    name=cuerda al aire,
    description={Pulsado de una cuerda sin presionar ningun traste}
}

\newglossaryentry{traste}
{
    name=traste,
    description={Tira metálica perpendicuar al mástil de un instrumento de cuerda pulsada}
}

\newglossaryentry{tablatura}
{
    name=tablatura,
    description={Notación musical para la representación gráfica a través de una tabla cuyas lineas horizontales equivalen a cuerdas y los números a los trastes}
}

\newglossaryentry{diagrama de acordes}
{
    name=diagrama de acordes,
    description={Notación musical para la representación de acordes sobre texto acompañados de un diagrama en forma de tabla cuyas lineas verticales equivalen a cuerdas y horizontales a trastes}
}

\newglossaryentry{cejilla}{
    name=cejilla,
    description={Utilizacion de un mismo dedo para la pulsación de varias cuerdas en un mismo traste}
}

\newglossaryentry{mock}
{
    name=mock,
    description={Un \textit{mock} es un objeto simulado que imita el comportamiento de un componente externo, útil para la ejecución de pruebas donde se pretende simular componentes externos.}
}